% !PRINTED 147
\weberpage{147}
%<<CONT

\weq{25}{\vartheta_{00}\left(0,\frac{\omega}{2}\right)^{2}=\vartheta_{00}^{2}+\vartheta_{10}^{2}=(1+\varkappa)\,\vartheta_{00}^{2}.}

Nach \S\,42 erhalten wir also eine Beziehung zwischen elliptischen Funktionen mit zwei verschiedenen Moduln, und es ist, wenn $\lambda$, $L$, $L'$ aus $\varkappa$, $K$, $K'$ durch die Vertauschung $\left(\omega,\dfrac{\omega}{2}\right)$ hervorgehen, nach (23), (24), (25)

\weq{26}{\lambda=\frac{2\sqrt{\varkappa}}{1+\varkappa},\quad L=(1+\varkappa)\,K,\quad L'=(1+\varkappa)\,K'}

\weq{27}{\operatorname{sn}\bigl[(1+\varkappa)\,v,\ \lambda\bigr]=\frac{(1+\varkappa)\operatorname{sn}v}{1+\varkappa\operatorname{sn}^{2}v},}

wenn der Deutlichkeit halber der Modul unter dem Funktionszeichen sn als zweites Argument mit aufgeführt ist.

Für die \gesperrt{Landen}sche Transformation ist
\[
\begin{aligned}
\frac{\vartheta_{11}(2u,\,2\omega)}{\vartheta_{01}(2u,\,2\omega)}&=\frac{\vartheta_{10}(u)\,\vartheta_{11}(u)}{\vartheta_{00}(u)\,\vartheta_{01}(u)},\\[4pt]
\frac{\vartheta_{01}(0,\,2\omega)^{2}}{\vartheta_{00}(0,\,2\omega)^{2}}&=\frac{2\,\vartheta_{00}\vartheta_{01}}{\vartheta_{00}^{2}+\vartheta_{01}^{2}}=\frac{2\sqrt{\varkappa'}}{1+\varkappa'},\\[4pt]
2\,\vartheta_{00}^{2}(0,\,2\omega)&=(1+\varkappa')\,\vartheta_{00}^{2},
\end{aligned}
\]
daraus, wenn $\lambda$, $\lambda'$, $L$, $L'$ durch die Vertauschung $(\omega,2\omega)$ aus $\varkappa$, $\varkappa'$, $K$, $K'$ entstehen:

\weq{28}{\begin{aligned}
\lambda'&=\frac{2\sqrt{\varkappa'}}{1+\varkappa'}, \qquad \sqrt{\lambda}=\frac{\varkappa}{1+\varkappa'}, \qquad \lambda=\frac{1-\varkappa'}{1+\varkappa'},\\[4pt]
2L&=(1+\varkappa')\,K, \qquad\qquad\qquad L'=(1+\varkappa')\,K',
\end{aligned}}

\weq{29}{\operatorname{sn}\bigl[(1+\varkappa')\,v,\ \lambda\bigr]=\frac{(1+\varkappa')\operatorname{sn}v\operatorname{cn}v}{\operatorname{dn}v},}

und in (26) bis (29) erkennt man nach \S\,42 die Formeln (2) und (4) des \S\,9.

\begin{center}
\textbf{\S\,45. \gesperrt{Die lineare Transformation der elliptischen Funktionen.}}
\end{center}

Die Einwirkung der linearen Transformation auf die elliptischen Funktionen übersieht man am besten aus der Darstellung durch die $\sigma$-Funktionen. Man erhält aus (4) oder (10), \S\,42 und den Formeln (4), (6) des \S\,37:
%CONT>>
